%% kapa: exam=1 round=35 year=2024 period=2 form=A subject=accounting subject_ko=회계학 seq=28 number=68
%% source: Q-Net 감정평가사 기출문제 (cst00309) artlSeq=5213229 file=raw/1008/5213229/2교시 문제지 원본.pdf
%% conversion: kordoc 4.24.51 -> tools/gen1.py (auto); review=proofread-sonnet; flags=none
\begin{question}{68}
매각예정으로 분류된 비유동자산 또는 처분자산집단의 회계처리에 관한 설명으로 옳지 않은 것은?
\begin{choices}
  \item 매각예정으로 분류된 비유동자산(또는 처분자산집단)은 공정가치에서 처분부대원가를 뺀 금액과 장부금액 중 큰 금액으로 측정한다.
  \item 1년 이후에 매각될 것으로 예상된다면 처분부대원가는 현재가치로 측정하고, 기간 경과에 따라 발생하는 처분부대원가 현재가치의 증가분은 금융원가로서 당기손익으로 회계처리한다.
  \item 매각예정으로 분류하였으나 중단영업의 정의를 충족하지 않는 비유동자산(또는 처분자산집단)을 재측정하여 인식하는 평가손익은 계속영업손익에 포함한다.
  \item 비유동자산이 매각예정으로 분류되거나 매각예정으로 분류된 처분자산집단의 일부이면 그 자산은 감가상각(또는 상각)하지 아니한다.
  \item 매각예정으로 분류된 처분자산집단의 부채와 관련된 이자와 기타 비용은 계속해서 인식한다.
\end{choices}
\end{question}
\answer{1}
